
\subsubsection{6.1 Complete Orthogonality}

The Jaccard similarity of 0.0 indicates that static analysis and execution feedback detect entirely non-overlapping error classes in this analysis. Static tools analyze code structure to identify type mismatches, undefined variables, and style violations. Execution tests evaluate runtime behavior to identify incorrect outputs and exceptions. These mechanisms operate on fundamentally different code properties, which explains the observed orthogonality.

The finding that 48.9\% of problems have static-only errors while 0\% have execution-only errors reflects that all execution failures in this dataset co-occur with static errors. This is consistent with the high static coverage (98.6\%).

\subsubsection{6.2 Static Analysis Reliability}

Static analysis detected errors in 98.6\% of code that failed execution tests. This high coverage suggests that static analysis is a reliable first-pass filter for identifying problematic code. The per-benchmark breakdown shows consistent performance: 100\% coverage on HumanEval+ and 98.0\% on MBPP+.

\subsubsection{6.3 Behavioral Detection Not Validated}

The behavioral detection rate of 0.3\% is far below the 40\% threshold. This result does not invalidate the hypothesis that execution detects behavioral errors in static-clean code; rather, it reflects a limitation of the experimental design.

\textbf{Root cause:} Canonical solutions from EvalPlus were analyzed instead of actual LLM-generated code. Canonical solutions are specifically designed to pass the extended test suites, so behavioral failures are rare by construction.

\textbf{Implication:} Validating H-M2 requires code with natural semantic errors---that is, actual LLM-generated code where behavioral errors can occur independently of structural errors.

\subsubsection{6.4 Limitations}

\textbf{L1: Canonical solutions used.} The use of canonical solutions rather than LLM-generated code prevents validation of the behavioral detection hypothesis. Future work should use code generated via LLM API calls.

\textbf{L2: Combined improvement untested.} The hypothesis that combined feedback exceeds the maximum of individual feedback types was not empirically tested. The orthogonality finding provides theoretical support but not empirical validation.

\textbf{L3: Single benchmark family.} Results are limited to Python code on HumanEval+ and MBPP+. Generalization to other languages and benchmarks is not established.

\textbf{L4: Single code sample per problem.} Analysis used one canonical solution per problem. Results may differ with multiple samples per problem.

